\section{Discussion and Conclusion}
\label{sec:conclusion}

We asked what a driving world model trained on raw video knows about the future
ego trajectory, using a small generative probe on its frozen features.

The collapse of the predicted distribution, which motivated this work, turned
out to have two separate causes. Sampling with too few ODE steps pulls every
sample towards the conditional mean, and that alone made a well-calibrated probe
look as if it was collapsing. Real collapse is memorization of the training
scenes. It needs a probe that can fit its training set, and its onset grows
linearly with the number of training scenes, as the theory of Bonnaire et
al.~\cite{bonnaire} predicts. A polynomial trajectory code is a better target
than raw poses but does not change this, and reducing the probe size only
prevents collapse by underfitting. The remedies are more data, early stopping on
a proper scoring rule, and enough sampling steps. Mean displacement error is
blind to all of this, so it should not be used to select generative trajectory
models.

The Orbis world model holds more about the future trajectory than DINOv2,
V-JEPA~2 and its own tokenizer, by a wide margin and at every depth from block 7
on. The information is concentrated in the middle of the network and in the
tokens of the last observed frame rather than in the frame being predicted.
Linear read-out confirms that ego speed and acceleration are almost directly
available there. A planner on these frozen features, trained for two minutes,
reaches an EPDMS of 0.881 on navtest, in the range of planners trained end to
end with more sensors.

Two results point to what is still missing. First, both the probe and the
planner produce good candidate trajectories (best of four or eight scores 0.90
to 0.96) but cannot pick the best one by imitation. Learning the simulator's
sub-scores is the known fix, and our attempt with 5\% of the targets was not
enough. Second, the probes do not generalize to unseen combinations of city and
command, and they stay confident when they fail. Their dispersion drops on such
inputs, so a calibration check on held-out data could at least flag the problem,
but the model itself gives no warning.

\paragraph{Limitations.}
Most probes were trained with one seed, so differences below a few percent in
energy are not reliable. The fast planner's three seeds span 0.02 EPDMS. The
memorization slope rests on three training-set sizes. All encoders were reduced
to 64 principal components, which may cost the wide encoders more than the
narrow ones. We use only the front camera, our navhard numbers cannot be
compared with the leaderboard, and our navtest numbers are a few points higher
than the published devkit gives. Finally, NAVSIM is open loop with replayed
agents, so none of this shows how the planners would behave in closed loop.
